\subsection{Covectors and the 1-form}
\lecdate{Lec 17 - Mar 12 (Week 9)}
... are we reviewing?

to understand VF, need to understand them as:
\begin{itemize}
	\item \textit{coord-wise}: $X=\sum X_{i}\p_{i}$, where $X_{i}$'s are smooth fns
		(in any chart on $M$)
	\item \textit{abstract}: a smooth section of the tangent bundle
		(smooth choice of tangent vec for each point)
	\item \textit{functional} [or sheaf-theoretic?]: a derivation on $C^{\infty}(M)$
	\item \textit{integral}: sth that generates a flow
		(a 1-param family of diffeos)
\end{itemize}

we can do sth similar for diff forms.
\begin{itemize}
	\item \textit{coord-wise}: $k$-form $\omega=\sum f_{I}\di\omega_{i_{1}}\wedge \cdots \wedge \d\omega_{i_{k}}$
	\item \textit{abstract}: $\omega$ is a smooth choice of $k$-covec at each point
		(a section of some [the cotangent] vector bundle)
	\item \textit{functional}: $\omega$ a kind of fn $(\frk{X}(M))^{k}\gto C^{\infty}(M)$
	\item \textit{integral}: a $k$-form is sth integrable over a $k$-mfld
\end{itemize}

\begin{xmp}
	given $f\in C^{\infty}(M)$, defn its (exterior) differential
	$\d f:\frk{X}(M)\gto C^{\infty}(M)$ by:
	\begin{equation*}
		\d f(X) = X(f)
	\end{equation*}
	note, $\d f$ is $C^{\infty}(M)$-linear, i.e.:
	\begin{equation*}
		\d f(gX+hY)=g \d f(X)+h\d f(Y)
	\end{equation*}
	that is an exercise.
\end{xmp} \

\begin{defn}[title=Functional]
	a \textbf{1-form} on $M$ is a $C^{\infty}(M)$-linear map $\frk{X}(M)\gto C^{\infty}(M)$.
	the set ($C^{\infty}(M)$-alg) of 1-forms is denoted $\Omega^{1}(M)$.
\end{defn}
hey, rmbr dual vspaces?
given $\alpha\in V^{*}$, we can write $\alpha(v)$ as $\ang{\alpha,v}$.
this emphasizes that if $V$ finite-dim, $(V^{*})^{*}=V$ canonically.
that is, we can think of $V$ as a hom $V^{*}\gto\bR$.
note $V^{*}\simeq V$, but not canonically unless we first fix an inner prod.

\begin{defn}
	given basis $e_{1},\dots,e_{n}$ of $V$, the \textbf{dual basis} for $V^{*}$
	is $e^{1},\dots,e^{n}$, where $e^{i}(e_{j})=\delta_{ij}$.
\end{defn}
cotangent space is defnd precisely as u think.
its elems are called \textbf{covectors}.
in $T_{p}\bR^{n}$, we have a distinguished basis $\p_{1},\dots,\p_{n}$.
its dual basis is written $\d x_{1},\dots,\d x_{n}$, where:
\begin{equation*}
	\d x_{i}(a_{1}\p_{1}+\dots +a_{n}\p_{n})=a_{i}
\end{equation*}

\begin{defn}
	given a lin. map $F:V\gto W$, theres a \textbf{dual map} $F^{*}:W^{*}\gto V^{*}$
	defn'd by $F^{*}\omega(v)=\omega(Fv)$.
	[compare this to pullback defn/notation]
\end{defn}
given smooth map $f:M\gto N$, we have a linear map $T_{p}f:T_{p}M\gto T_{f(p)}N$
given by $v\mto f_{*}v=T_{p}f(v)$.
therefore, $T_{p}^{*}f:T_{f(p)}^{*}N\gto T_{p}^{*}M$ is given as
$v\mto f^{*}v=T_{p}^{*}f(v)$.
the ``base" map is as a pushforward, i.e. vecs are \textit{covariant}.
the dual is as a pullback, i.e. covecs are \textit{contravariant}.

\begin{xmp}
	given $f\in C^{\infty}(M)$, we can defn a covec $\d_{p}f\in T^{*}_{p}M$ by:
	\begin{equation*}
		\d_{p}f(v)=v(f)=D_{v}f
	\end{equation*}
	in particular, in $\bR^{n}$, we can write
	$\d_{p}f=\p_{x_{1}}f\d x_{1}+\dots +\p_{x_{n}}f\d x_{n}$.
	exercise... lol.
\end{xmp} \

\begin{prop}
	given $f\in C^{\infty}(M,N),g\in C^{\infty}(N)$, we have
	$\d_{p}(f^{*}g)=(T_{p}^{*}f)(\d_{p}g)$.
\end{prop} \

\begin{pf}
	for any $v\in T_{p}M$:
	\begin{align*}
		\ang{\d_{p}(f^{*}g),v}
		& \ = \ v(f^{*}g) \\
		& \ = \ v(g\circ f) \\
		& \ = \ (T_{p}f(v))(g) \\
		& \ = \ \ang{\d_{f(p)}g,T_{p}f(v)} \\
		& \ = \ \ang{(T_{p}^{*}f)(\d_{f(p)}g),v}
	\end{align*}
	note $\ang{\cdot,\cdot}$ here is evaluation, not an inner prod.
\end{pf}

\lecdate{Lec 18 - Mar 17 (Week 10)}
last time, we talked abt covecs - this is enough for a ``vibe" of a defn.
what is a 1-form?
it is a smooth choice of covec at each pt of a mfld.

\begin{defn}[title=Coordinate-wise]
	a map $p\mto\alpha_{p}$ sending pts to covecs is a \textbf{diff 1-form} if
	for every chart $(U,\vphi)$:
	\begin{equation*}
		\alpha_{p} \ = \ \sum_{i=1}^{m}g_{i}(\vphi(p))\di x_{i}
	\end{equation*}
	where $g_{i}:\vphi(U)\gto\bR$ is smooth.
\end{defn}

for instance:
\begin{enumerate}[1)]
	\item $f:M\gto\bR$ as $\d f = \sum\p_{x_{i}}f-\d x_{i}$ on a chart.
		as an exercise, show this is indep of choice of chart.
	\item $y\d x$ on $\bR^{2}$ is not $\d f$ for any $f$.
		to see this, note its partials don't commute [why?], or that
		$y\d x$ depends on $y$ so $f$ depends on $y$, so $\p_{y}f\neq0$.
		[what?]
\end{enumerate}

\begin{defn}[title=Functional]
	a 1-form is a $C^{\infty}(M)$-linear $\frk{X}(M)\gto C^{\infty}(M)$
	given by applying the covec to pts.
\end{defn} \

\begin{prop}
	let $\alpha$ be a 1-form in the functional sense.
	then for any VF $X$, there exists a well-defn'd covec $\alpha_{p}$ w $p\in M$ s.t.
	$\alpha(X)(p)=\alpha_{p}(X_{p})$.
\end{prop} \

\begin{pf}
	wts if $X_{p}=Y_{p}$, then $\alpha(X)(p)=\alpha(Y)(p)$.
	therefore, this gives a well-defn'd map $X_{p}\mto\alpha_{p}(X_{p})$.
	that it's linear, ie a covec, is an exercise.
	by linearity, enough to show that if $X_{p}=0$ then $\alpha(X)(p)=0$.

	indeed, if $X_{p}=0$, write $X=\sum f_{i}Y_{i}$ where $f_{i}(p)=0$
	($Y_{i}=\p_{i}$ on a chart near $p$, then extend somehow).
	thus:
	\begin{equation*}
		\alpha(X)(p) \ = \ \sum_{i}f_{i}(p)\alpha(Y_{i})(p) \ = \ 0
	\end{equation*}
\end{pf} \

\begin{prop}\vspace{-0.2in}
	\begin{enumerate}[a)]
		\item functional 1-forms can be restricted to a form on an open
			set $U\subseteq M$.
		\item the two defns are equiv.
	\end{enumerate}
\end{prop}
note that this is non-trivial since a VF on $U$ is not always a restriction of
a VF on $M$.
\begin{exr}
	state carefully what the above means, then prove it.
\end{exr}

given $g:N\gto\bR$, recall a pullback of $g$ along $f:M\gto N$ is defined as:
\begin{equation*}
	f^{*}g \ = \ g\circ f
\end{equation*}
for covecs, there is an induced map $T_{p}f:T_{p}M\gto T_{f(p)}N$ by
$v\mto f_{*}v$, which in turn induces a map $T_{f(p)}^{*}N\gto T_{p}^{*}M$ by
$\nu\mto f^{*}\nu$, where $\nu$ is a covec.

now, given a 1-form $\omega\in\Omega^{1}(N)$, we can define its \textbf{pullback} along $f$ by:
\begin{equation*}
	(f^{*}\omega)_{p} \ = \ f^{*}(\omega_{f(p)})
\end{equation*}
we claim this is indeed a 1-form, i.e. its a smooth choice.
indeed, in coords, we have a smooth fn $f:U\gto V$ where
$U,V\subseteq\bR^{m},\bR^{n}$ resp.
if $\omega\in\Omega^{1}(V)$ given by $\sum g_{i}\d x_{i}$, then $T_{p}f$ has smooth
coord fns $\p_{x_{j}}f_{i}$. note:
\begin{equation*}
	f^{*}\big\rvert_{f(p)} \ = \ (T_{p}f)^{\intercal}
\end{equation*}
so:
\begin{equation*}
	(f^{*}\omega)_{p} \ = \ f^{*}\bigg\rvert_{f(p)}(\omega_{f(p)})
	\ = \ f^{*}\bigg\rvert_{f(p)}\sum g_{i}(f(p))\di x_{i}
\end{equation*}
these are all smooth in $p$, so we are done.

now, some properties of pullbacks:
\begin{itemize}
	\item $(f\circ g)^{*} \ = \ g^{*}\circ f^{*}$
	\item if $g:N\gto\bR$ smooth, $f^{*}(g\omega) \ = \ f^{*}g \cdot f^{*}\omega$
	\item $f^{*}(\d g) \ = \ \d(f^{*}g)$
	\item in coords:
		\begin{align*}
			(f^{*}\omega)_{p} \ = \ f^{*}\bigg\rvert_{f(p)}\sum_{i}g_{i}(f(p))\di x_{i}
			& \ = \ \sum_{i}g_{i}(f(p))f^{*}\bigg\rvert_{f(p)}\di x_{i} \\
			& \ = \ \sum_{i}\sum_{j}g_{i}(f(p))\frac{\p f_{i}}{\p x_{j}}\di x_{j}
		\end{align*}
\end{itemize}

\begin{defn}
	the \textbf{cotangent bundle} of a mfld $M$ is a mfld structure on:
	\begin{equation*}
		T^{*}M \ := \ \bigsqcup_{p\in M}T_{p}^{*}M
	\end{equation*}
	with the canonical proj $\pi:T^{*}M\gto M$.
	given $(U,\vphi)$, we get a chart on $T^{*}M$:
	\begin{equation*}
		(\pi^{-1}(U),(\vphi^{-1})^{*}:\pi^{-1}U\gto\vphi(U)\times\bR^{m})
	\end{equation*}
	here, $\bR^{m}$ is a copy of the cotangent space.
	note transition maps are smooth, as they are given by:
	\begin{equation*}
		(\vphi\circ\psi^{-1})^{*} = (T_{p}(\vphi\circ\psi^{-1}))^{\intercal} \ : \ \vphi(U\cap V)\times\bR^{m}\gto\psi(U\cap V)\times\bR^{m}
	\end{equation*}
\end{defn}
this is just repackaging the coord-wise defn.

NOTE: with tangent bundles, $f:M\gto N$ defined $Tf:TM\gto TN$.
here, we can only do sth similar when $f$ is a diffeo, as pts go forward but
vecs go back.
OTOH, if $\omega\in\Omega^{1}(N)$ is defn'd as $s_{\omega}:N\gto T^{*}N$, then
$f^{*}\omega$ is given by:
\begin{equation*}
	s_{f^{*}\omega}(p) \ = \ f^{*}\big\rvert_{f(p)}(s_{\omega}\circ f(p))
\end{equation*}

\begin{defn}[title=Abstract]
	a 1-form is a section of the cotangent bundle
\end{defn}

given $\gamma:I\gto M$, $\alpha\in\Omega^{1}(M)$, take $\gamma^{*}\alpha\in\Omega^{1}(I)$ as
$\gamma^{*}(\alpha)=f(t)\di t$ for some $f(t)$.

with this, we can define the integral of a 1-form.

\begin{defn}
	we define:
	\begin{equation*}
		\int_{\gamma}\alpha \ = \ \int_{I}\gamma^{*}\alpha \ = \ \int_{a}^{b}f(t)\di t
	\end{equation*}
	where $I=[a,b]$.
	we'll take smoothness on closed intervals in the extension way.
\end{defn} \

\begin{prop}
	if $\alpha=\d f$ for some $f:M\gto\bR$, then:
	\begin{equation*}
		\int_{\gamma}\alpha \ = \ f(\gamma(b))-f(\gamma(a))
	\end{equation*}
\end{prop} \

\begin{pf}
	we have that
	$\gamma^{*}\d f \ = \ \d(\gamma^{*}f) \ = \ \d(f\circ\gamma) \ = \ \p_{t}(f\circ\gamma)\d t$.
	apply usual ftc.
\end{pf}

for instance, $y\d x$ is not $\d f$, i.e. an \textit{exact form}, since
by integrating counterclockwise on $\p[0,1]^{2}$, we get:
\begin{equation*}
	\int_{\gamma}y\di x \ = \ 0+0+\int_{0}^{1}1(-\d t) + 0 \ = \ 1 \ \neq \ 0 \ = \ f(0,0)-f(0,0)
\end{equation*}

\begin{prop}
	integrating is invariant wrt reparametrization:
	\begin{equation*}
		\int_{\gamma}\alpha \ = \ \pm\int_{\gamma\circ\kappa}\alpha
	\end{equation*}
	where $\kappa:[c,d]\gto[a,b]$ is a diffeo, and $\pm$ depends on orientation.
\end{prop} \

\begin{pf}
	we have:
	\begin{equation*}
		\int_{\gamma\circ\kappa}\alpha \ = \ \int_{c}^{d}(\gamma\circ\kappa)^{*}\alpha
		\ = \ \int_{c}^{d}\kappa^{*}\gamma^{*}\alpha \ = \ \pm \int_{\kappa(c)}^{\kappa(d)}\gamma^{*}\alpha
	\end{equation*}
	after a substitution.
\end{pf}

recall $\alpha$ is \textit{exact} if its $\d f$ for some $f\in C^{\infty}$.
in this case, $\int_{\gamma}\alpha$ depends only on endpts.

\textit{fact: this is an iff}.
if $\int_{\gamma}\alpha$ depends only on endpts, then defn $f(p)=\int_{\gamma}\alpha$ where
$\gamma$ is any curve $p_{0}$ to $p$ (where we can choose $p_{0}$).
to ensure $f$ smooth and $\d f=\alpha$, do a coord calculation.
ofc, this is iff $\int_{\gamma}\alpha=0$ for $\gamma$ a loop.

\begin{xmp}
	let $M=S^{1},\alpha=\d\theta$. then:
	\begin{equation*}
		\int_{S^{1}}\d\theta \ = \ \int_{\gamma}\d\theta \ = \ 2\pi
	\end{equation*}
	where $\gamma$ inj on $(a,b)$ and $\gamma(a)=\gamma(b)$.
	thus, $\d\theta$ is not an exact form.
\end{xmp}

